\section{Experiments}
We demonstrate \modelname is an effective architecture for  universal image segmentation  through comparisons with specialized state-of-the-art architectures on  standard benchmarks.
We evaluate %
our proposed design decisions through ablations on all three tasks.
Finally we show \modelname generalizes beyond the standard benchmarks,
obtaining state-of-the-art results on four datasets.

\begin{table*}[t]
  \centering

  \tablestyle{4pt}{1.2}\scriptsize\begin{tabular}{l | lcc | x{20}x{20}x{20} | x{20}x{24} |x{24}x{24}x{18}}
  method & backbone & query type & epochs & PQ & PQ$^\text{Th}$ & PQ$^\text{St}$ & AP$^\text{Th}_\text{pan}$ & mIoU$_\text{pan}$ & \#params. & FLOPs & fps \\
  \shline
  \demph{DETR~\cite{detr}} & \demph{R50} & \demph{100 queries} & \demph{500+25} & \demph{43.4} & \demph{48.2} & \demph{36.3} & \demph{31.1} & \demph{-} & \demph{-} & \demph{-} & \demph{-} \\
  MaskFormer~\cite{cheng2021maskformer} & R50 & 100 queries & 300 & 46.5 & 51.0 & 39.8 & 33.0 & 57.8 & \phantom{0}45M & \phantom{0}181G & 17.6 \\
  \textbf{\modelname} (ours) & R50 & 100 queries & 50 & \textbf{51.9} & \textbf{57.7} & \textbf{43.0} & \textbf{41.7} & \textbf{61.7} & \phantom{0}44M & \phantom{0}226G & \phantom{0}8.6 \\
  \hline
  \demph{DETR~\cite{detr}} & \demph{R101} & \demph{100 queries} & \demph{500+25} & \demph{45.1} & \demph{50.5} & \demph{37.0} & \demph{33.0} & \demph{-} & \demph{-} & \demph{-} & \demph{-} \\
  MaskFormer~\cite{cheng2021maskformer} & R101 & 100 queries & 300 & 47.6 & 52.5 & 40.3 & 34.1 & 59.3 & \phantom{0}64M & \phantom{0}248G & 14.0 \\
  \textbf{\modelname} (ours) & R101 & 100 queries & 50 & \textbf{52.6} & \textbf{58.5} & \textbf{43.7} & \textbf{42.6} & \textbf{62.4} & \phantom{0}63M & \phantom{0}293G & \phantom{0}7.2 \\
  \hline
  Max-DeepLab~\cite{wang2021max} & Max-L & 128 queries & 216 & 51.1 & 57.0 & 42.2 & - & - & 451M & 3692G & \demph{-} \\
  MaskFormer~\cite{cheng2021maskformer} & Swin-L$^{\text{\textdagger}}$ & 100 queries & 300 & 52.7 & 58.5 & 44.0 & 40.1 & 64.8 & 212M & \phantom{0}792G & \phantom{0}5.2 \\
  K-Net~\cite{zhang2021knet} & Swin-L$^{\text{\textdagger}}$ & 100 queries & 36 & 54.6 & 60.2 & 46.0 & - & - & - & - & - \\
  \textbf{\modelname} (ours) & Swin-L$^{\text{\textdagger}}$ & \emph{200 queries} & \emph{100} & \textbf{57.8} & \textbf{64.2} & \textbf{48.1} & \textbf{48.6} & \textbf{67.4} & 216M & \phantom{0}868G & \phantom{0}4.0 \\
  \end{tabular}
  \vspace{-1mm}

   \caption{\textbf{Panoptic segmentation on COCO panoptic \texttt{val2017} with 133 categories.} \modelname consistently outperforms MaskFormer~\cite{cheng2021maskformer} by a large margin with different backbones on all metrics. Our best model outperforms prior state-of-the-art MaskFormer by $5.1$ PQ and K-Net~\cite{zhang2021knet} by $3.2$ PQ.
Backbones pre-trained on ImageNet-22K are marked with $^{\text{\textdagger}}$.
}
\vspace{-3mm}

\label{tab:panseg:coco}
\end{table*}

\noindent\textbf{Datasets.}
We study \modelname %
using four widely used image segmentation datasets that support semantic, instance and panoptic segmentation:
COCO~\cite{lin2014coco} (80 ``things'' and 53 ``stuff'' categories), ADE20K~\cite{zhou2017ade20k} (100 ``things'' and 50 ``stuff'' categories), Cityscapes~\cite{Cordts2016Cityscapes} (8 ``things'' and 11 ``stuff'' categories) and Mapillary Vistas~\cite{neuhold2017mapillary} (37 ``things'' and 28 ``stuff'' categories). Panoptic and semantic segmentation tasks are evaluated on the union of ``things'' and ``stuff'' categories while instance segmentation is only evaluated on the ``things'' categories.

\noindent\textbf{Evaluation metrics.}
For \emph{panoptic segmentation}, we use the standard \textbf{PQ} (panoptic quality) metric~\cite{kirillov2017panoptic}. We further report \textbf{AP$^\text{Th}_\text{pan}$}, which is the AP evaluated on the ``thing'' categories using instance segmentation annotations, and \textbf{mIoU$_\text{pan}$}, which is the mIoU for semantic segmentation by merging instance masks from the same category, of the same model trained \emph{only} with panoptic segmentation annotations. For \emph{instance segmentation}, we use the standard \textbf{AP} (average precision) metric~\cite{lin2014coco}. For \emph{semantic segmentation}, we use \textbf{mIoU} (mean Intersection-over-Union)~\cite{everingham2015pascal}.


\begin{table*}[t]
  \centering

  \tablestyle{4pt}{1.2}\scriptsize\begin{tabular}{l| lcc | x{20}x{20}x{20}x{20} | x{30} |x{24}x{24}x{18}}
  method & backbone & query type & epochs & AP & AP$^\text{S}$ & AP$^\text{M}$ & AP$^\text{L}$ & AP$^\text{boundary}$ & \#params. & FLOPs & fps \\
  \shline
  MaskFormer~\cite{cheng2021maskformer} & R50 & 100 queries & 300 & 34.0 & 16.4 & 37.8 & 54.2 & 23.0 & \phantom{0}45M & \phantom{0}181G & 19.2 \\
  \demph{Mask R-CNN~\cite{he2017mask}} & \demph{R50} & \demph{dense anchors} & \demph{36} & \demph{37.2} & \demph{18.6} & \demph{39.5} & \demph{53.3} & \demph{23.1} & \demph{\phantom{0}44M} & \demph{\phantom{0}201G} & \demph{15.2} \\
  Mask R-CNN~\cite{he2017mask,ghiasi2021simple,du2021simple} & R50 & dense anchors & 400 & 42.5 & \textbf{23.8} & 45.0 & 60.0 & 28.0 & \phantom{0}46M & \phantom{0}358G & 10.3 \\
  \textbf{\modelname} (ours) & R50 & 100 queries & 50 & \textbf{43.7} & 23.4 & \textbf{47.2} & \textbf{64.8} & \textbf{30.6} & \phantom{0}44M & \phantom{0}226G & \phantom{0}9.7 \\
  \hline
  \demph{Mask R-CNN~\cite{he2017mask}} & \demph{R101} & \demph{dense anchors} & \demph{36} & \demph{38.6} & \demph{19.5} & \demph{41.3} & \demph{55.3} & \demph{24.5} & \demph{\phantom{0}63M} & \demph{\phantom{0}266G} & \demph{10.8} \\
  Mask R-CNN~\cite{he2017mask,ghiasi2021simple,du2021simple} & R101 & dense anchors & 400 & 43.7 & \textbf{24.6} & 46.4 & 61.8 & 29.1 & \phantom{0}65M & \phantom{0}423G & \phantom{0}8.6 \\
  \textbf{\modelname} (ours) & R101 & 100 queries & 50 & \textbf{44.2} & 23.8 & \textbf{47.7} & \textbf{66.7} & \textbf{31.1} & \phantom{0}63M & \phantom{0}293G & \phantom{0}7.8 \\
  \hline
  QueryInst~\cite{QueryInst} & Swin-L$^{\text{\textdagger}}$ & 300 queries & 50 & 48.9 & 30.8 & 52.6 & 68.3 & 33.5 & - & - & \phantom{0}\demph{3.3} \\
  Swin-HTC++~\cite{liu2021swin,chen2019hybrid} & Swin-L$^{\text{\textdagger}}$ & dense anchors & 72 & 49.5 & \textbf{31.0} & 52.4 & 67.2 & 34.1 & 284M & 1470G & - \\
  \textbf{\modelname} (ours) & Swin-L$^{\text{\textdagger}}$ & \emph{200 queries} & \emph{100} & \textbf{50.1} & 29.9 & \textbf{53.9} & \textbf{72.1} & \textbf{36.2} & 216M & \phantom{0}868G & \phantom{0}4.0 \\
  \end{tabular}
  \vspace{-1mm}

   \caption{\textbf{Instance segmentation on COCO \texttt{val2017} with 80 categories.} \modelname outperforms strong Mask R-CNN~\cite{he2017mask} baselines for both AP and AP$^\text{boundary}$~\cite{cheng2021boundary} metrics when training with $8\x$ fewer epochs. Our best model is also competitive to the state-of-the-art specialized instance segmentation model on COCO and has higher boundary quality. For a fair comparison, we only consider single-scale inference and models trained using only COCO \texttt{train2017} set data.
Backbones pre-trained on ImageNet-22K are marked with $^{\text{\textdagger}}$.}

\vspace{-3mm}

\label{tab:insseg:coco}
\end{table*}





\subsection{Implementation details}
\label{sec:exp:impl}
\noindent We adopt settings from~\cite{cheng2021maskformer} with the following differences:


\noindent\textbf{Pixel decoder.} \modelname is compatible with %
any existing pixel decoder module. %
In MaskFormer~\cite{cheng2021maskformer}, FPN~\cite{lin2016feature} is chosen as the default %
for its simplicity. Since our goal is to demonstrate strong performance across different segmentation tasks, %
we use the more advanced multi-scale deformable attention Transformer (MSDeformAttn)~\cite{zhu2021deformable} as our default pixel decoder. Specifically, we use 6 MSDeformAttn layers applied to feature maps with resolution $1/8$, $1/16$ and $1/32$,
and use a simple upsampling layer with lateral connection on the final %
$1/8$ feature map to generate the  feature map of resolution $1/4$ as the per-pixel embedding.
In our ablation study, we show that this pixel decoder provides best results across different segmentation tasks.

\noindent\textbf{Transformer decoder.} We use our Transformer decoder proposed in \secref{sec:method:arch:transformer_decoder} with $L=3$ (\ie, 9 layers total) and 100 queries by default. An auxiliary loss is added to every intermediate Transformer decoder layer and to the learnable query features before the Transformer decoder.

\noindent\textbf{Loss weights.} We use the binary cross-entropy loss (instead of focal loss~\cite{lin2017focal} in~\cite{cheng2021maskformer}) and the dice loss~\cite{milletari2016v} for our mask loss: $\mathcal{L}_{\text{mask}} = \lambda_{\text{ce}}\mathcal{L}_{\text{ce}} +\lambda_{\text{dice}}\mathcal{L}_{\text{dice}}$. We set $\lambda_{\text{ce}} = 5.0$ and $\lambda_{\text{dice}} = 5.0$. The final loss is a combination of mask loss and classification loss: $\mathcal{L}_{\text{mask}} + \lambda_{\text{cls}}\mathcal{L}_{\text{cls}}$ and we set $\lambda_{\text{cls}} = 2.0$ for predictions matched with a ground truth and $0.1$ for the ``no object,'' \ie, predictions that have not been matched with any ground truth.


\noindent\textbf{Post-processing.} We use the exact same post-processing as~\cite{cheng2021maskformer} to acquire the expected output format for panoptic and semantic segmentation from pairs of binary masks and class predictions. Instance segmentation requires additional confidence scores for each prediction. We multiply class confidence and mask confidence (\ie, averaged foreground per-pixel binary mask probability) for a final confidence. %

\subsection{Training settings}
\noindent\textbf{Panoptic and instance segmentation.} We use Detectron2~\cite{wu2019detectron2} and follow the updated Mask R-CNN~\cite{he2017mask} baseline settings\footnote{\scriptsize \url{https://github.com/facebookresearch/detectron2/blob/main/MODEL\_ZOO.md\#new-baselines-using-large-scale-jitter-and-longer-training-schedule}} for the COCO dataset. More specifically, we use AdamW~\cite{loshchilov2018decoupled} optimizer and the step learning rate  schedule. We use an initial learning rate of $0.0001$ and a weight decay of $0.05$ for all backbones. A learning rate multiplier of $0.1$ is applied to the backbone and we decay the learning rate at 0.9 and 0.95 fractions of the total number of training steps by a factor of 10. If not stated otherwise, we train our models for 50 epochs with a batch size of 16. For data augmentation, we use the large-scale jittering (LSJ) augmentation~\cite{ghiasi2021simple,du2021simple} with a random scale sampled from range 0.1 to 2.0 followed by a fixed size crop to $1024\x1024$. We use the standard Mask R-CNN inference setting where we resize an image with shorter side to 800 and longer side up-to 1333. We also report FLOPs and fps. FLOPs are averaged over 100 validation images (COCO images have varying sizes). Frames-per-second (fps) is measured on a V100 GPU with a batch size of 1 by taking the average runtime on the entire validation set including post-processing time.

\noindent\textbf{Semantic segmentation.} We follow the same settings as~\cite{cheng2021maskformer} to train our models, except: 1) a learning rate multiplier of 0.1 is applied to \emph{both} CNN and Transformer backbones instead of only applying it to CNN backbones in~\cite{cheng2021maskformer},
2) both ResNet and Swin backbones use an initial learning rate of $0.0001$ and a weight decay of $0.05$, instead of using different learning rates in~\cite{cheng2021maskformer}.

\subsection{Main results}

\noindent\textbf{Panoptic segmentation.} We compare \modelname with state-of-the-art models for panoptic segmentation on the COCO panoptic~\cite{kirillov2017panoptic} dataset in \tabref{tab:panseg:coco}. \modelname consistently outperforms MaskFormer by more than 5 PQ across different backbones while converging $6\x$ faster. With Swin-L backbone,
our \modelname sets a new state-of-the-art of 57.8 PQ, outperforming existing state-of-the-art~\cite{cheng2021maskformer} by 5.1 PQ and concurrent work, K-Net~\cite{zhang2021knet}, by 3.2 PQ.
\modelname even outperforms the best ensemble models with extra training data in the COCO challenge (see \appref{app:results:panoptic} for test set results).

Beyond the PQ metric, our \modelname also achieves higher performance on two other metrics compared to DETR~\cite{detr} and MaskFormer: AP$^\text{Th}_\text{pan}$, which is the AP evaluated on the 80 ``thing'' categories using \emph{instance segmentation annotation}, and mIoU$_\text{pan}$, which is the mIoU evaluated on the 133 categories for semantic segmentation converted from panoptic segmentation annotation. This shows \modelname's universality: %
trained \emph{only} with panoptic segmentation annotations, it can be used for instance and semantic segmentation.

\begin{table}[t]
  \centering
  \tablestyle{2pt}{1.2}
  \scriptsize
  \begin{tabular}{l | lc | x{35}x{35}}
  method & backbone & crop size & mIoU (s.s.) & mIoU (m.s.) \\
  \shline
  MaskFormer~\cite{cheng2021maskformer} & R50 & $512$ & 44.5 & 46.7 \\
  \textbf{\modelname} (ours) & R50 & $512$ & \textbf{47.2} & \textbf{49.2} \\
  \hline\hline
  \demph{Swin-UperNet~\cite{liu2021swin,xiao2018unified}} & \demph{Swin-T\phantom{$^{\text{\textdagger}}$}} & \demph{$512$} & \demph{-} & \demph{46.1} \\
  MaskFormer~\cite{cheng2021maskformer} & Swin-T\phantom{$^{\text{\textdagger}}$} & $512$ & 46.7 & 48.8 \\
  \textbf{\modelname} (ours) & Swin-T & $512$ & \textbf{47.7} & \textbf{49.6} \\
  \hline\hline
  MaskFormer~\cite{cheng2021maskformer} & Swin-L$^{\text{\textdagger}}$ & $640$ & 54.1 & 55.6 \\
  FaPN-MaskFormer~\cite{fapn,cheng2021maskformer} & Swin-L-FaPN$^{\text{\textdagger}}$ & $640$ & 55.2 & 56.7 \\
  BEiT-UperNet~\cite{beit,xiao2018unified} & BEiT-L$^{\text{\textdagger}}$ & $640$ & - & 57.0 \\
  \hline
  \multirow{2}{*}{\textbf{\modelname} (ours)}
  & Swin-L$^{\text{\textdagger}}$ & $640$ & 56.1 & 57.3 \\
  & Swin-L-FaPN$^{\text{\textdagger}}$ & $640$ & \textbf{56.4} & \textbf{57.7} \\
  \end{tabular}

  \caption{\textbf{Semantic segmentation on ADE20K \texttt{val} with 150 categories.} \modelname consistently outperforms MaskFormer~\cite{cheng2021maskformer} by a large margin with different backbones (all \modelname models use MSDeformAttn~\cite{zhu2021deformable} as pixel decoder, except Swin-L-FaPN uses FaPN~\cite{fapn}). Our best model outperforms the best specialized model, BEiT~\cite{beit}. We report both single-scale (s.s.) and multi-scale (m.s.) inference results.
  Backbones pre-trained on ImageNet-22K are marked with $^{\text{\textdagger}}$.}
  \vspace{-3mm}

\label{tab:semseg:ade20k}
\end{table}


\begin{table*}[t]
  \begin{subtable}{0.49\linewidth}
  \centering
  \tablestyle{3pt}{1.2}
  \scriptsize
  \begin{tabular}{y{80} | x{30}x{30}x{30} | x{30}}
   & AP & PQ & mIoU & FLOPs \\
  \shline
  \textbf{\modelname} (ours) & \textbf{43.7} \phantom{\dt{-0.0}} & \textbf{51.9} \phantom{\dt{-0.0}} & \textbf{47.2} \phantom{\dt{-0.0}} & 226G \\
  \hline
  $-$ masked attention & 37.8 \dt{-5.9} & 47.1 \dt{-4.8} & 45.5 \dt{-1.7} & 213G \\
  $-$ high-resolution features & 41.5 \dt{-2.2} & 50.2 \dt{-1.7} & 46.1 \dt{-1.1} & 218G \\
  \multicolumn{5}{c}{~}\\
  \multicolumn{5}{c}{~}\\
  \end{tabular}
  \caption{Masked attention and high-resolution features (from efficient multi-scale strategy) lead to the most gains. More detailed ablations are in \tabref{tab:ablation:maskformer:a} and \tabref{tab:ablation:maskformer:b}. We remove one component at a time.
  }
  \label{tab:ablation:transformer:a}
  \end{subtable}\hspace{2mm}
  \begin{subtable}{0.49\linewidth}
  \centering
  \tablestyle{3pt}{1.2}
  \scriptsize
  \begin{tabular}{y{80} | x{30}x{30}x{30} | x{30}}
  & AP & PQ & mIoU & FLOPs \\
  \shline
  \modelname (ours) & \textbf{43.7} \phantom{\dt{-0.0}} & \textbf{51.9} \phantom{\dt{-0.0}} & \textbf{47.2} \phantom{\dt{-0.0}} & 226G \\
  \hline
  $-$ learnable query features & 42.9 \dt{-0.8} & 51.2 \dt{-0.7} & 45.4 \dt{-1.8} & 226G \\
  $-$ cross-attention first & 43.2 \dt{-0.5} & 51.6 \dt{-0.3} & 46.3 \dt{-0.9} & 226G \\
  $-$ remove dropout & 43.0 \dt{-0.7} & 51.3 \dt{-0.6} & 47.2 \dt{-0.0} & 226G \\
  \hline
  $-$ all 3 components above & 42.3 \dt{-1.4} & 50.8 \dt{-1.1} & 46.3 \dt{-0.9} & 226G \\
  \end{tabular}
  \caption{Optimization improvements increase the performance without introducing extra compute. Following DETR~\cite{detr}, query features are zero-initialized when not learnable. We remove one component at a time.}
  \label{tab:ablation:transformer:b}
  \end{subtable}\vspace{2mm}
  \begin{subtable}{0.3\linewidth}
  \centering
  \tablestyle{1pt}{1.2}
  \scriptsize
  \begin{tabular}{l | x{20}x{20}x{20} | x{24}}
   & AP & PQ & mIoU & FLOPs \\
  \shline
  cross-attention & 37.8 & 47.1 & 45.5 & 213G \\
  \hline
  SMCA~\cite{gao2021smca} & 37.9 & 47.2 & \underline{46.6} & 213G \\
  mask pooling~\cite{zhang2021knet} & \underline{43.1} & \underline{51.5} & 46.0 & 217G \\
  \hline
  \textbf{masked attention} & \textbf{43.7} & \textbf{51.9} & \textbf{47.2} & 226G \\
  \multicolumn{5}{c}{~}\\
  \end{tabular}
  \caption{\textbf{Masked attention.} Our masked attention performs better than other variants of cross-attention across all tasks.
  }
  \label{tab:ablation:maskformer:a}
  \end{subtable}\hspace{2mm}
  \begin{subtable}{0.35\linewidth}
  \centering
  \tablestyle{1pt}{1.2}
  \scriptsize
  \begin{tabular}{l | x{20}x{20}x{20} | x{24}}
   & AP & PQ & mIoU & FLOPs \\
  \shline
  single scale ($1/32$) & 41.5 & 50.2 & 46.1 & 218G \\
  single scale ($1/16$) & 43.0 & 51.5 & 46.5 & 222G \\
  single scale ($1/8$) & \textbf{44.0} & \underline{51.8} & \textbf{47.4} & 239G \\
  \hline
  na\"ive m.s.\ (3 scales) & \textbf{44.0} & \textbf{51.9} & 46.3 & 247G \\
  \hline
  \textbf{efficient m.s.} (3 scales) & \underline{43.7} & \textbf{51.9} & \underline{47.2} & 226G \\
  \end{tabular}
  \caption{\textbf{Feature resolution.}
  High-resolution features (single scale $1/8$) are important. Our efficient multi-scale (efficient m.s.) strategy effectively reduces the FLOPs.
  }
  \label{tab:ablation:maskformer:b}
  \end{subtable}\hspace{2mm}
  \begin{subtable}{0.32\linewidth}
  \centering
  \tablestyle{1pt}{1.2}
  \scriptsize
  \begin{tabular}{l | x{20}x{20}x{20} | x{24}}
   & AP & PQ & mIoU & FLOPs \\
  \shline
  FPN~\cite{lin2016feature} & 41.5 & 50.7 & 45.6 & 195G \\
  \hline
  Semantic FPN~\cite{kirillov2019panopticfpn} & 42.1 & 51.2 & 46.2 & 258G \\
  FaPN~\cite{fapn} & 42.4 & \underline{51.8} & \underline{46.8} & - \\
  BiFPN~\cite{tan2020efficientdet} & \underline{43.5} & \underline{51.8} & 45.6 & 204G \\
  \hline
  \textbf{MSDeformAttn}~\cite{zhu2021deformable} & \textbf{43.7} & \textbf{51.9} & \textbf{47.2} & 226G \\
  \end{tabular}
  \caption{\textbf{Pixel decoder.} MSDeformAttn~\cite{zhu2021deformable} consistently performs the best across all tasks.
  \phantom{pad} \phantom{pad} \phantom{pad} \phantom{pad} \phantom{pad} \phantom{pad} \phantom{pad} \phantom{pad} \phantom{pad} %
  }
  \label{tab:ablation:maskformer:c}
  \end{subtable}
  \caption{\textbf{\modelname ablations.} We perform ablations on three tasks: instance (AP on COCO \texttt{val2017}), panoptic (PQ on COCO panoptic \texttt{val2017}) and semantic (mIoU on ADE20K \texttt{val}) segmentation. FLOPs are measured on COCO instance segmentation.
  }
  \label{tab:ablation:maskformer}
  \vspace{-3mm}
\end{table*}




\noindent\textbf{Instance segmentation.} We compare \modelname with state-of-the-art models %
on the COCO~\cite{lin2014coco} dataset in \tabref{tab:insseg:coco}. With ResNet~\cite{he2016deep} backbone, \modelname outperforms a strong Mask R-CNN~\cite{he2017mask} baseline using large-scale jittering (LSJ) augmentation~\cite{ghiasi2021simple,du2021simple} while requiring $8\x$ fewer training iterations.
With Swin-L backbone, \modelname outperforms the state-of-the-art HTC++~\cite{chen2019hybrid}.
Although we only observe $+0.6$ AP improvement over HTC++, the Boundary AP~\cite{cheng2021boundary} improves by $2.1$, suggesting that our predictions have a better boundary quality %
thanks to the high-resolution mask predictions.
Note that for a fair comparison, we only consider single-scale inference and models trained with only COCO \texttt{train2017} set data.

With a ResNet-50 backbone \modelname  improves over MaskFormer on small objects by 7.0 AP$^\text{S}$, while overall the highest gains come from large objects ($+10.6$ AP$^\text{L}$). The performance on AP$^\text{S}$ still lags behind other state-of-the-art models.
Hence there still remains room for improvement on small objects, \eg, by using dilated backbones like in DETR~\cite{detr}, which we leave for future work.

\noindent\textbf{Semantic segmentation.} We compare \modelname with state-of-the-art models for semantic segmentation on the ADE20K~\cite{zhou2017ade20k} dataset in \tabref{tab:semseg:ade20k}. \modelname outperforms MaskFormer~\cite{cheng2021maskformer} across different backbones, suggesting that the proposed improvements even boost semantic segmentation results where~\cite{cheng2021maskformer} was already state-of-the-art. With Swin-L as  backbone and FaPN~\cite{fapn} as  pixel decoder, \modelname sets a new state-of-the-art of 57.7 mIoU. We also report the test set results in \appref{app:results:semantic}.


\subsection{Ablation studies}
We now analyze \modelname through a series of ablation studies using a ResNet-50 backbone~\cite{he2016deep}. To test the generality of the proposed components  for universal image segmentation, \emph{all ablations are performed on three tasks}.





\noindent\textbf{Transformer decoder.}
We validate the importance of each component %
by removing them one at a time. As shown in \tabref{tab:ablation:transformer:a}, masked attention leads to the biggest improvement across all tasks. The improvement is larger for %
instance and panoptic segmentation
than for semantic segmentation. Moreover, using high-resolution features from the efficient multi-scale strategy is also important. \tabref{tab:ablation:transformer:b} shows additional optimization improvements further improve the performance without extra computation.

\noindent\textbf{Masked attention.}
Concurrent work has proposed other variants of cross-attention~\cite{gao2021smca,meng2021cdetr} that aim to improve the convergence and performance of DETR~\cite{detr} for object detection. Most recently, K-Net~\cite{zhang2021knet} replaced cross-attention with a mask pooling operation that averages features within mask regions. We validate the importance of our masked attention in \tabref{tab:ablation:maskformer:a}. While existing cross-attention variants may improve on a specific task, our masked attention performs the best on all three tasks.


\noindent\textbf{Feature resolution.}
\tabref{tab:ablation:maskformer:b} shows that \modelname benefits from using high-resolution features (\eg, a single scale of $1/8$) in the Transformer decoder. However, this introduces additional computation. Our efficient multi-scale (efficient m.s.) strategy effectively reduces the FLOPs without affecting the performance. Note that, naively concatenating multi-scale features as input to every Transformer decoder layer (na\"ive m.s.) does not yield additional gains.

\noindent\textbf{Pixel decoder.}
As shown in \tabref{tab:ablation:maskformer:c}, \modelname is compatible with any existing pixel decoder. However, we observe different pixel decoders specialize in different tasks: while BiFPN~\cite{tan2020efficientdet} performs better on instance-level segmentation, FaPN~\cite{fapn} works better for semantic segmentation. Among all studied pixel decoders, the MSDeformaAttn~\cite{zhu2021deformable} consistently performs the best across all tasks and thus is selected as our default. This set of ablations also suggests that designing a module like a pixel decoder for a specific task does not guarantee generalization across  segmentation tasks. \modelname, as a universal model, could serve as a testbed for a generalizable module design.

\begin{table}[t]
  \centering
  \tablestyle{3pt}{1.2}
  \scriptsize
  \begin{tabular}{x{40}x{40}| x{30}x{30}x{30} | x{24}}
  \phantom{matching loss} matching loss & \phantom{matching loss} training loss & AP (COCO) & PQ (COCO) & mIoU (ADE20K) & memory (COCO) \\
  \shline
  \multirow{2}{*}{mask \phantom{(ours)}} & mask \phantom{(ours)} & 41.0 & 50.3 & 45.9 & 18G \\ %
   & point \phantom{(ours)} & 41.0 & 50.8 & 45.9 & \phantom{0}6G \\ %
  \hline
  \multirow{2}{*}{\textbf{point} (ours)} & mask \phantom{(ours)} & 43.1 & 51.4 & \textbf{47.3} & 18G \\ %
   & \textbf{point} (ours) & \textbf{43.7} & \textbf{51.9} & 47.2 & \phantom{0}6G \\ %
  \end{tabular}
  \caption{\textbf{Calculating loss with points \vs masks.} Training with point loss reduces training memory without influencing the performance. Matching with point loss further improves performance.}
  \label{tab:ablation:matching}
  \vspace{-3mm}
\end{table}


\noindent\textbf{Calculating loss with points \vs masks.}
In \tabref{tab:ablation:matching} we study the performance and memory implications when calculating the loss based on either mask or sampled points. Calculating the final training loss with sampled points reduces training memory by $3\x$ without affecting the performance. Additionally, calculating the matching loss with sampled points improves performance across all three tasks.


\begin{figure}[t]
    \begin{subtable}{1.0\linewidth}
    \centering
    \bgroup
    \def\arraystretch{0.2}
    \setlength\tabcolsep{0.2pt}
    \begin{tabular}{cccc}
    \includegraphics[width=0.24\linewidth]{figures/vis_query_before_transformer/img_000000127517/query_6.jpg} &
    \includegraphics[width=0.24\linewidth]{figures/vis_query_before_transformer/img_000000127517/query_20.jpg} &
    \includegraphics[width=0.24\linewidth]{figures/vis_query_before_transformer/img_000000127517/query_73.jpg} &
    \includegraphics[width=0.24\linewidth]{figures/vis_query_before_transformer/img_000000127517/query_97.jpg} \\
    \includegraphics[width=0.24\linewidth]{figures/vis_query_before_transformer/img_000000218424/query_6.jpg} &
    \includegraphics[width=0.24\linewidth]{figures/vis_query_before_transformer/img_000000218424/query_20.jpg} &
    \includegraphics[width=0.24\linewidth]{figures/vis_query_before_transformer/img_000000218424/query_73.jpg} &
    \includegraphics[width=0.24\linewidth]{figures/vis_query_before_transformer/img_000000218424/query_97.jpg} \\
    \includegraphics[width=0.24\linewidth]{figures/vis_query_before_transformer/img_000000377814/query_6.jpg} &
    \includegraphics[width=0.24\linewidth]{figures/vis_query_before_transformer/img_000000377814/query_20.jpg} &
    \includegraphics[width=0.24\linewidth]{figures/vis_query_before_transformer/img_000000377814/query_73.jpg} &
    \includegraphics[width=0.24\linewidth]{figures/vis_query_before_transformer/img_000000377814/query_97.jpg} \\
    \end{tabular} \egroup
    \end{subtable}\vspace{0mm}
    \begin{subtable}{0.48\linewidth}
      \centering
      \tablestyle{3pt}{1.2}
      \scriptsize
      \begin{tabular}{c}
       \includegraphics[width=0.99\linewidth]{figures/layer_wise_ar.pdf} \\
      \end{tabular}
      \caption*{AR@100 on COCO \texttt{val2017}}
    \end{subtable}
    \begin{subtable}{0.48\linewidth}
      \centering
      \tablestyle{3pt}{1.2}
      \scriptsize
      \begin{tabular}{c}
       \includegraphics[width=0.99\linewidth]{figures/maskformerv2_proposals.pdf} \\
      \end{tabular}
    \end{subtable}
  \vspace{-2mm}
  \caption{
  \textbf{Learnable queries as ``region proposals''.} \emph{Top:} We visualize mask predictions of four selected learnable queries \emph{before} feeding them into the Transformer decoder (using R50 backbone).
  \emph{Bottom left:} We calculate the class-agnostic average recall with 100 proposals (AR@100) and observe that these learnable queries provide good proposals compared to the final predictions of \modelname after the Transformer decoder layers (layer 9).
  \emph{Bottom right:} Illustration of proposal generation process.
  }
  \label{fig:analysis:query}
\end{figure}




\noindent\textbf{Learnable queries as region proposals.} Region proposals~\cite{uijlings2013selective,arbelaez2014multiscale}, either in the form of boxes or masks, are regions that are likely to be ``objects.'' With learnable queries being supervised by the mask loss, predictions from learnable queries can serve as mask proposals. In \figref{fig:analysis:query} top, we visualize mask predictions of selected learnable queries \emph{before} feeding them into the Transformer decoder (the proposal generation process is shown in \figref{fig:analysis:query} bottom right). In \figref{fig:analysis:query} bottom left, we further perform a quantitative analysis on the quality of these proposals by calculating the class-agnostic average recall with 100 predictions (AR@100) on COCO \texttt{val2017}. We find these learnable queries already achieve good AR@100 compared to the final predictions of \modelname after the Transformer decoder layers, \ie, layer 9, and AR@100 consistently improves with more decoder layers.


\begin{table}[t]
  \centering

  \tablestyle{1pt}{1.2}\scriptsize\begin{tabular}{l|l | x{16}x{20}x{24} |x{36}x{16}}
  & & \multicolumn{3}{c|}{panoptic model} & \multicolumn{2}{c}{semantic model} \\
  method & backbone & PQ & AP$^\text{Th}_\text{pan}$ & mIoU$_\text{pan}$ & mIoU (s.s.) & (m.s.) \\
  \shline
  Panoptic FCN~\cite{li2021fully} & Swin-L$^{\text{\textdagger}}$ & 65.9\phantom{$^*$} & - & - & - & - \\
  Panoptic-DeepLab~\cite{cheng2020panoptic} & SWideRNet~\cite{chen2020scaling} & 66.4\phantom{$^*$} & 40.1\phantom{$^*$} & 82.2\phantom{$^*$} & - & - \\
  \demph{Panoptic-DeepLab~\cite{cheng2020panoptic}} & \demph{SWideRNet~\cite{chen2020scaling}} & \demph{67.5$^*$} & \demph{43.9$^*$} & \demph{82.9$^*$} & - & - \\
  \hline
  SETR~\cite{zheng2021rethinking} & ViT-L$^{\text{\textdagger}}$~\cite{dosovitskiy2020vit} & - & - & - & - & 82.2 \\
  SegFormer~\cite{xie2021segformer} & MiT-B5~\cite{xie2021segformer} & - & - & - & - & 84.0 \\
  \hline\hline
  \multirow{3}{*}{\textbf{\modelname} (ours)}
  & R50\phantom{$^{\text{\textdagger}}$} & 62.1\phantom{$^*$} & 37.3\phantom{$^*$} & 77.5\phantom{$^*$} & 79.4 & 82.2 \\
  & Swin-B$^{\text{\textdagger}}$ & 66.1\phantom{$^*$} & 42.8\phantom{$^*$} & 82.7\phantom{$^*$} & \textbf{83.3} & \textbf{84.5} \\
  & Swin-L$^{\text{\textdagger}}$ & \textbf{66.6}\phantom{$^*$} & \textbf{43.6}\phantom{$^*$} & \textbf{82.9}\phantom{$^*$} & \textbf{83.3} & 84.3 \\
  \end{tabular}
  \vspace{-2mm}

   \caption{\textbf{Cityscapes \texttt{val}.} \modelname is competitive to specialized models on Cityscapes. Panoptic segmentation models use single-scale inference by default, multi-scale numbers are marked with~$^*$. For semantic segmentation, we report both single-scale (s.s.) and multi-scale (m.s.) inference results. Backbones pre-trained on ImageNet-22K are marked with $^{\text{\textdagger}}$.}
  \vspace{-3mm}

\label{tab:benchmark:cityscapes}
\end{table}

\subsection{Generalization to other datasets}
To show our \modelname can generalize beyond the COCO dataset, we further perform experiments on other popular image segmentation datasets. In \tabref{tab:benchmark:cityscapes}, we show results on Cityscapes~\cite{Cordts2016Cityscapes}. Please see \appref{app:datasets} for detailed training settings on each dataset as well as more results on  ADE20K~\cite{zhou2017ade20k} and Mapillary Vistas~\cite{neuhold2017mapillary}.

We observe that our \modelname is competitive to state-of-the-art methods on these datasets as well. It suggests \modelname can serve as a universal image segmentation model and results generalize across datasets.

\begin{table}[t]
\centering
  \begin{subtable}{0.4\linewidth}
  \centering
  \tablestyle{2pt}{1.2}
  \scriptsize
  \begin{tabular}{l | x{16}x{16}x{16}}
   & PQ & AP & mIoU \\
  \shline
  panoptic & \demph{51.9} & 41.7 & \textbf{61.7} \\
  \hline
  instance & - & \textbf{43.7} & - \\
  semantic & - & - & 61.5 \\
  \end{tabular}
  \caption{COCO}
  \label{tab:analysis:universal:coco}
  \end{subtable}
  \begin{subtable}{0.25\linewidth}
  \centering
  \tablestyle{1pt}{1.2}
  \scriptsize
  \begin{tabular}{x{16}x{16}x{16}}
  PQ & AP & mIoU \\
  \shline
  \demph{39.7} & \textbf{26.5} & 46.1 \\
  \hline
  - & 26.4 & - \\
  - & - & \textbf{47.2} \\
  \end{tabular}
  \caption{ADE20K}
  \label{tab:analysis:universal:ade20k}
  \end{subtable}
  \begin{subtable}{0.25\linewidth}
  \centering
  \tablestyle{1pt}{1.2}
  \scriptsize
  \begin{tabular}{x{16}x{16}x{16}}
  PQ & AP & mIoU \\
  \shline
  \demph{62.1} & 37.3 & 77.5 \\
  \hline
  - & \textbf{37.4} & - \\
  - & - & \textbf{79.4} \\
  \end{tabular}
  \caption{Cityscapes}
  \label{tab:analysis:universal:cs}
  \end{subtable}
  \vspace{-2mm}
  \caption{\textbf{Limitations of \modelname.} Although a single \modelname can address any segmentation task, we still need to train it on different tasks. Across three datasets we find \modelname trained with panoptic annotations performs slightly worse than the exact same model trained specifically for instance and semantic segmentation tasks with the corresponding data.}
  \label{tab:analysis:universal}
  \vspace{-3mm}
\end{table}



\subsection{Limitations}
Our ultimate goal is to train a {\em single} model for \emph{all} image segmentation tasks. In \tabref{tab:analysis:universal}, we find \modelname trained on panoptic segmentation only performs slightly worse than the exact same model trained with the corresponding annotations for instance and semantic segmentation tasks across three datasets.
This suggests that even though \modelname can generalize to different tasks, it still needs to be trained for those specific tasks.
In the future, we hope to develop a model that can be trained only once for multiple tasks and even for multiple datasets.

Furthermore, as seen in Tables~\ref{tab:insseg:coco} and \ref{tab:ablation:maskformer:b}, even though it improves over baselines, \modelname struggles with segmenting small objects and is unable to fully leverage multi-scale features. We believe better utilization of the feature pyramid and designing losses for small objects are critical.



















